\documentclass[11pt,a4paper]{article}

\usepackage[utf8]{inputenc}
\usepackage[T1]{fontenc}
\usepackage{amsmath,amssymb,amsfonts,amsthm}
\usepackage{booktabs}
\usepackage{float}
\usepackage{microtype}
\usepackage{tabularx}
\usepackage{geometry}
\usepackage{algorithm}
\usepackage{algpseudocode}
\usepackage{cite}
\usepackage{url}
\usepackage{listings}
\usepackage[dvipsnames]{xcolor}
\usepackage{hyperref}

\geometry{a4paper, margin=2.5cm}
\hypersetup{colorlinks=true, linkcolor=MidnightBlue, citecolor=MidnightBlue, urlcolor=MidnightBlue}

\theoremstyle{definition}
\newtheorem{definition}{Definition}
\newtheorem{theorem}{Theorem}
\newtheorem{proposition}{Proposition}
\newtheorem{remark}{Remark}

\newcommand{\Z}{\mathbb{Z}}
\newcommand{\Zq}{\mathbb{Z}_q}
\newcommand{\CBD}{\operatorname{CBD}}
\newcommand{\SHAKE}{\operatorname{SHAKE256}}
\newcommand{\SHA}{\operatorname{SHA256}}

\title{\textbf{Windowed Hankel KEM: A Compact Lattice-Based KEM from Truncated Linear Convolution}}
\author{Marcos Mendoza III\\ \small Independent Researcher}
\date{October 2026}

\begin{document}
\maketitle

\begin{abstract}
We present \emph{Windowed Hankel KEM}, a key encapsulation mechanism based on a structured variant of the Learning with Errors (LWE) problem. The public matrix is a Hankel matrix generated by a single pseudorandom sequence. Equivalently, the public key is a noisy \emph{middle product}: the fully supported coefficients of a linear (non-cyclic) convolution of the public sequence with a short secret.

This structure lies between unstructured LWE, as used by FrodoKEM, and ring- or module-LWE, as used by ML-KEM. It compresses the public key without reducing products modulo a fixed polynomial. Middle-Product LWE (MP-LWE), introduced by Ro\c{s}ca, Sakzad, Steinfeld and Stehl\'e, admits reductions from Polynomial-LWE over a large family of rings. However, our instance does not satisfy the hypotheses of those reductions: it uses a single long public sequence, a short binomial secret, and a narrow error distribution. We therefore claim no reduction, and we assess security by concrete lattice-attack estimates.

With $n = 640$, a window of $W = 128$ coordinates, $q = 2^{12}$, and binomial noise ($\eta = 2$), the scheme encapsulates a 256-bit message by 4-PAM modulation. Its public key is 1{,}196 bytes and its ciphertext 1{,}356 bytes, both including 12-byte format headers. A Chernoff bound gives a decryption failure rate below $2^{-134.1}$. Treating the instance as unstructured LWE, the cheapest attack found by the lattice estimator is the dual hybrid, with BKZ block size $\beta = 468$ and an estimated cost of $2^{137.1}$ in the Core-SVP model. Under the same model and tool, Kyber-512 costs $2^{115.5}$. A Python implementation with ARMv8 NEON kernels completes a handshake in 0.34~ms on an Apple M3 Max.
\end{abstract}

\paragraph{Keywords:} post-quantum cryptography, key encapsulation, learning with errors, middle-product LWE, Hankel matrices.

%=====================================================================
\section{Introduction}
%=====================================================================

Lattice-based key encapsulation mechanisms (KEMs) occupy a spectrum of algebraic structure.

At one end, unstructured LWE~\cite{regev2009lattices} uses a uniformly random public matrix $A \in \Zq^{m \times n}$ and is backed by worst-case reductions from standard lattice problems. FrodoKEM~\cite{bos2016frodo} follows this approach. Its Category~1 parameter set transmits 19{,}336 bytes per exchange (public key plus ciphertext).

At the other end, ring- and module-LWE schemes such as ML-KEM~\cite{fips203,bos2018kyber} replace $A$ by matrices of polynomials in $\Zq[X]/(X^{n}+1)$. A single ring element then describes an $n \times n$ block. This reduces ML-KEM-512 to 1{,}568 bytes per exchange and enables NTT-based multiplication. The resulting algebraic structure has been studied extensively; see~\cite{peikert2016decade} for a survey.

This paper studies an intermediate structure. Here the public matrix is a \emph{Hankel} matrix $H_{i,j} = a_{i+j}$ generated by a single sequence $a$. Multiplying by $H$ computes a contiguous slice of the coefficients of the polynomial product $a(X)\,s(X)$ in $\Zq[X]$, without reduction modulo any polynomial.

The choice of slice matters. The outermost coefficients of a full linear convolution depend on only a few secret coordinates; for instance, $c_0 = a_0 s_0 + e_0$. Publishing them would let an attacker recover a small secret coordinate by coordinate. We therefore publish only the \emph{fully supported} coefficients, each of which depends on all $n$ secret coordinates. This is precisely the \emph{middle product} of Ro\c{s}ca, Sakzad, Steinfeld and Stehl\'e~\cite{rsss2017middle}, who introduced Middle-Product LWE (MP-LWE). Steinfeld, Sakzad and Zhao later designed the Titanium cryptosystem around it~\cite{titanium2017nist,ssz2019practical}.

The MP-LWE reductions cover a specific sample distribution: many samples, each with a fresh short public polynomial; a long uniform secret; and Gaussian noise of sufficient width. Titanium uses parameters compatible with these reductions; its 128-bit KEM uses a 16{,}352-byte public key and a 3{,}552-byte ciphertext~\cite{titanium2017nist}. Our design deliberately departs from that distribution in order to be compact. It uses a single long public sequence, a short binomial secret, and narrow binomial noise. Its security is therefore a heuristic assumption, supported by concrete cryptanalysis rather than by a reduction.

\paragraph{Contributions.}
\begin{enumerate}
  \item A KEM based on windowed Hankel (middle-product) products, with a 128-coordinate window that carries a 256-bit message by 4-PAM modulation, and a modulus $q = 2^{12}$ that reduces all arithmetic to masking (Section~\ref{sec:scheme}).
  \item A precise account of the scheme's relation to MP-LWE, including why the known reductions do not apply to the chosen parameters (Section~\ref{sec:mplwe}).
  \item A decryption-failure analysis with an exact Chernoff bound, and concrete estimates for primal, dual and hybrid lattice attacks (Section~\ref{sec:security}).
  \item An IND-CCA KEM via the Fujisaki--Okamoto transform with implicit rejection~\cite{hhk2017modular}, and an ARMv8 NEON implementation (Section~\ref{sec:impl}).
\end{enumerate}

%=====================================================================
\section{Preliminaries}
%=====================================================================

\subsection{Hankel operators and the middle product}
Throughout, $q = 4096$ and $\Zq = \Z/q\Z$.

\begin{definition}[Hankel matrix]
A matrix $H \in \Zq^{m \times n}$ is a \emph{Hankel matrix} if $H_{i,j} = a_{i+j}$ for some sequence $a = (a_0, \dots, a_{m+n-2}) \in \Zq^{m+n-1}$.
\end{definition}

Write $a(X) = \sum_t a_t X^t$ and $s(X) = \sum_{j<n} s_j X^j$. Then $(Hs)_i = \sum_{j} a_{i+j} s_j$. This is the coefficient of $X^{i+n-1}$ in the product $a(X)\,\tilde{s}(X)$, where $\tilde{s}$ is the reversal of $s$. Hence $Hs$ is a contiguous slice of a polynomial product in $\Zq[X]$.

Choosing $a$ of length $m + n - 1$ ensures that every published coordinate depends on all $n$ entries of $s$. This is the \emph{middle product} of~\cite{rsss2017middle}. Hankel matrices are not closed under multiplication and do not form a ring. Nevertheless, the scheme is polynomial-structured: it uses polynomial multiplication in $\Zq[X]$, but without reduction modulo a fixed polynomial.

\subsection{Relation to Middle-Product LWE}\label{sec:mplwe}
In MP-LWE~\cite{rsss2017middle}, each sample consists of a fresh uniformly random short polynomial $a_i$ together with the middle product of $a_i$ and a fixed secret $s$ of roughly twice the length, plus Gaussian noise. Ro\c{s}ca et al.\ reduce Polynomial-LWE, over any polynomial in a large family, to MP-LWE. Consequently, MP-LWE is at least as hard as Polynomial-LWE for the hardest member of that family. Steinfeld et al.\ tighten this reduction and use it to set Titanium's parameters~\cite{ssz2019practical}.

The distribution used by Windowed Hankel KEM differs in three respects.
\begin{enumerate}
  \item A \emph{single} long public sequence $a$ of length $2n + W - 1$ is used, rather than many fresh short samples.
  \item The secret is \emph{short}, drawn from $\CBD_2$, rather than long and uniform.
  \item The error has standard deviation $1$, far below the width the reductions require.
\end{enumerate}
The reductions of~\cite{rsss2017middle,ssz2019practical} therefore do not apply, and we do not claim that the scheme inherits their guarantees. We regard the hardness of recovering $s$ from $(a, Hs + e)$, with $H$ windowed Hankel, as a heuristic assumption. Section~\ref{sec:security} evaluates it with concrete lattice attacks.

\subsection{Power-of-two modulus}
With $q = 2^{12}$, reduction modulo $q$ is a bitwise mask. Uniform elements are obtained by masking 16-bit words of $\SHAKE$~\cite{fips202} output, without rejection. Two 12-bit coefficients pack exactly into three bytes, so every byte string of the correct length decodes to canonical coefficients. The cost is that NTT-based multiplication is unavailable. Products are computed directly in $O(mn)$ operations (Section~\ref{sec:impl}).

\subsection{4-PAM modulation}
A 256-bit message is carried by $W = 128$ coordinates, two bits per coordinate. Symbol $\mu \in \{0,1,2,3\}$ maps to $\mu \cdot 1024 \in \Z_{4096}$.

\begin{proposition}[Decoder]
Let $w = 1024\mu + e \bmod 4096$ with $-512 \le e \le 511$. Then $\hat{\mu} = ((w + 512) \bmod 4096) \gg 10$ equals $\mu$.
\end{proposition}
\begin{proof}
We have $w + 512 \equiv 1024\mu + (e + 512) \pmod{4096}$ with $0 \le e + 512 \le 1023$. Hence $\lfloor ((w+512) \bmod 4096)/1024 \rfloor = \mu$, and division by $1024$ is a right shift by 10.
\end{proof}

%=====================================================================
\section{The Scheme}\label{sec:scheme}
%=====================================================================

The parameters are $n = 640$, $W = 128$, $m = n + W = 768$, $q = 4096$ and $\eta = 2$. The public sequence has length $L_a = 2n + W - 1 = 1407$. $H \in \Zq^{m \times n}$ denotes the Hankel matrix of $a$. $\CBD_2$ samples values in $\{-2,\dots,2\}$ with probabilities $(1,4,6,4,1)/16$, so its variance is $1$.

\begin{algorithm}[H]
\caption{CPA key generation}\label{alg:cpa_keygen}
\begin{algorithmic}[1]
\Require seed $\in \{0,1\}^{512}$
\State $\mathit{seed}_a \gets \mathit{seed}[0{:}32]$; \ $\mathit{seed}_e \gets \mathit{seed}[32{:}64]$
\State $a \gets (\SHAKE(\mathit{seed}_a, 2L_a)$ read as 16-bit words$) \bmod 4096$
\State $\beta \gets \SHAKE(\mathit{seed}_e \,\|\, \texttt{"\_hankel\_noise"}, n + m)$
\State $s \gets \CBD_2(\beta[0{:}n])$; \ $e \gets \CBD_2(\beta[n{:}n+m])$
\State $b \gets H s + e \bmod q$
\State \Return $\mathit{pk} = (\mathit{seed}_a, b)$, \ $\mathit{sk} = s$
\end{algorithmic}
\end{algorithm}

\begin{algorithm}[H]
\caption{CPA encryption of a 32-byte message $M$ with coins $c$}\label{alg:cpa_encaps}
\begin{algorithmic}[1]
\Require $\mathit{pk} = (\mathit{seed}_a, b)$, message $M$, coins $c \in \{0,1\}^{256}$
\State Regenerate $a$ from $\mathit{seed}_a$
\State $\gamma \gets \SHAKE(c \,\|\, \texttt{"\_hankel\_encaps\_coins"}, n + m + W)$
\State $r \gets \CBD_2(\gamma[0{:}n])$; \ $e_1 \gets \CBD_2(\gamma[n{:}n+m])$; \ $e_2 \gets \CBD_2(\gamma[n+m{:}n+m+W])$
\State $u \gets H r + e_1 \bmod q$
\State $K[w] \gets \sum_{j=0}^{n-1} r_j\, b_{w+j} \bmod q$ for $0 \le w < W$
\State $v \gets K + \mathrm{PAM}(M) + e_2 \bmod q$
\State \Return $(u, v)$
\end{algorithmic}
\end{algorithm}

\begin{algorithm}[H]
\caption{CPA decryption}\label{alg:cpa_decaps}
\begin{algorithmic}[1]
\Require $\mathit{sk} = s$, ciphertext $(u, v)$
\State $K'[w] \gets \sum_{k=0}^{n-1} s_k\, u_{w+k} \bmod q$ for $0 \le w < W$
\State \Return $\mathrm{Decode}(v - K' \bmod q)$
\end{algorithmic}
\end{algorithm}

\paragraph{Correctness.}
Expanding $b$ and $u$ gives
\[
K[w] = \sum_{j,k} a_{w+j+k}\, r_j s_k + \sum_j r_j e_{w+j},
\qquad
K'[w] = \sum_{j,k} a_{w+j+k}\, r_j s_k + \sum_k s_k e_{1,w+k}.
\]
The double sums coincide because each depends only on $j + k$. Therefore
\[
v - K' = \mathrm{PAM}(M) + \underbrace{\sum_{j=0}^{n-1} r_j e_{w+j} - \sum_{k=0}^{n-1} s_k e_{1,w+k} + e_{2,w}}_{e_{\mathrm{dec}}[w]} \pmod q,
\]
and decryption succeeds whenever $|e_{\mathrm{dec}}[w]| < 512$ for all $w$.

The structure $K \approx r^\top H s$ is the standard one for LWE-based key exchange~\cite{lindner2011better}. The term $e_2$ is essential. Without it, the $W$ coordinates of $v$ would give an attacker $W$ noiseless linear equations in the ephemeral secret $r$.

\paragraph{IND-CCA KEM.}
We apply the Fujisaki--Okamoto transform with implicit rejection, $\mathrm{QFO}_m^{\not\perp}$~\cite{hhk2017modular}.
\begin{itemize}
  \item \emph{Encapsulation} derives $(\bar{K}, c) \gets \SHAKE(M \,\|\, \SHA(\mathit{pk}), 64)$ for uniform $M$, and encrypts $M$ with coins $c$.
  \item \emph{Decapsulation} decrypts to $M'$ and re-encrypts. It compares the result with the received ciphertext in constant time, and selects either $\bar{K}'$ or a secret rejection value without branching.
  \item The shared secret is $\SHAKE(K_{\mathrm{sel}} \,\|\, \SHA(\mathit{ct}), 32)$.
\end{itemize}
Malformed ciphertexts are parsed without data-dependent exceptions: length, header and parameter checks are combined with bitwise operations.

%=====================================================================
\section{Parameters and Sizes}
%=====================================================================

\begin{table}[H]
\centering
\begin{tabular}{lc}
\toprule
Parameter & Value \\
\midrule
Dimension $n$ & 640 \\
Window $W$ / published coordinates $m = n + W$ & 128 / 768 \\
Modulus $q$ & $4096 = 2^{12}$ \\
Secret and error distribution & $\CBD_2$ (variance 1) \\
Public sequence length $2n + W - 1$ & 1{,}407 \\
Message & 256 bits (4-PAM over $W$ coordinates) \\
\midrule
Public key & 1{,}196 B \\
Ciphertext & 1{,}356 B \\
Secret key & 1{,}868 B \\
\bottomrule
\end{tabular}
\caption{Parameters of Windowed Hankel KEM. The public key consists of a 12-byte header, a 32-byte seed and $768$ coefficients of 12 bits; the ciphertext of a 12-byte header and $768 + 128$ coefficients of 12 bits.}
\label{tab:params}
\end{table}

Table~\ref{tab:comparison} compares wire sizes with LWE-based KEMs at the 128-bit level. For a like-for-like comparison, the Windowed Hankel figures in that table exclude our 12-byte format headers.

\begin{table}[H]
\centering
\small
\begin{tabular}{llrrr}
\toprule
Scheme & Structure & Public key & Ciphertext & Total \\
\midrule
ML-KEM-512~\cite{fips203} & module-LWE & 800 B & 768 B & 1{,}568 B \\
Windowed Hankel KEM (this work) & middle product & 1{,}184 B & 1{,}344 B & 2{,}528 B \\
Titanium-CCA-Std128~\cite{titanium2017nist} & MP-LWE & 16{,}352 B & 3{,}552 B & 19{,}904 B \\
FrodoKEM-640~\cite{bos2016frodo} & unstructured LWE & 9{,}616 B & 9{,}720 B & 19{,}336 B \\
\bottomrule
\end{tabular}
\caption{Wire sizes of LWE-based KEMs targeting the 128-bit security level.}
\label{tab:comparison}
\end{table}

The ciphertext is transmitted uncompressed. Rounding $u$ and $v$ to fewer bits, as ML-KEM does, would shrink it further but requires a revised failure analysis. We leave this to future work.

%=====================================================================
\section{Security Analysis}\label{sec:security}
%=====================================================================

\subsection{Lattice attacks}
Following the Core-SVP methodology~\cite{adps2016newhope,aps2015concrete}, we first treat the public key as an instance of unstructured LWE with $n = 640$, up to $m = 768$ samples, $q = 4096$, and secret and error of standard deviation $1$. The same estimate applies to the ciphertext component $u$, whose secret is $r$.

\paragraph{Estimates.}
We use the lattice estimator~\cite{latticeestimator} (commit \texttt{53da598}, August 2026) under two cost models for BKZ: the Core-SVP model of~\cite{adps2016newhope}, which charges $2^{0.292\beta}$ classically for one call to an SVP oracle in dimension $\beta$, and the more detailed model of MATZOV~\cite{matzov2022}. Table~\ref{tab:attacks} lists the results. For calibration, we ran the same version of the estimator on its built-in Kyber-512 parameters; the Kyber-512 figures it produces coincide with those in the estimator's documentation.

\begin{table}[H]
\centering
\small
\begin{tabular}{l rr rr rr}
\toprule
& \multicolumn{4}{c}{Windowed Hankel KEM} & \multicolumn{2}{c}{Kyber-512} \\
\cmidrule(lr){2-5}\cmidrule(lr){6-7}
& \multicolumn{2}{c}{Core-SVP} & \multicolumn{2}{c}{MATZOV} & Core-SVP & MATZOV \\
Attack & $\beta$ & cost & $\beta$ & cost & cost & cost \\
\midrule
Primal (uSVP)  & 488 & $2^{142.5}$ & 488 & $2^{166.8}$ & $2^{118.6}$ & $2^{143.8}$ \\
Primal (BDD)   & 488 & $2^{143.2}$ & 471 & $2^{163.1}$ & $2^{119.4}$ & $2^{140.2}$ \\
Dual           & 507 & $2^{148.0}$ & 509 & $2^{173.7}$ & $2^{123.8}$ & $2^{149.9}$ \\
Dual hybrid    & 468 & $2^{137.1}$ & 459 & $2^{160.3}$ & $2^{115.5}$ & $2^{139.7}$ \\
Coded-BKW      & --- & ---         & --- & $2^{207.3}$ & ---         & $2^{178.8}$ \\
\bottomrule
\end{tabular}
\caption{Lattice-estimator results for the public key ($n = 640$, $m = 768$, $q = 4096$, $\CBD_2$ secret and error) and for Kyber-512, treated as unstructured LWE. Costs are the estimator's total cost (\texttt{rop}), including guessing in the hybrid attack. Coded-BKW is costed without lattice reduction, so it appears only once per scheme.}
\label{tab:attacks}
\end{table}

The dual hybrid attack is the cheapest in both models. In the Core-SVP model it uses $\beta = 468$, so its lattice-reduction cost alone is $2^{0.292 \cdot 468} = 2^{136.7}$ classically and $2^{0.265 \cdot 468} = 2^{124.0}$ quantumly. The primal attack uses $\beta = 488$ in a lattice of dimension $1{,}194$, at a classical Core-SVP cost of $2^{142.5}$ and a quantum one of $2^{129.3}$. Kyber-512 reports a classical Core-SVP cost of $2^{118}$ for its Category~1 claim~\cite{bos2018kyber}, in agreement with the estimator's primal figure. In every row of Table~\ref{tab:attacks}, the Windowed Hankel instance costs at least 20 bits more than Kyber-512 under the same model; the smallest margin, 20.6 bits, is for the dual hybrid attack in the MATZOV model. These figures assume that the Hankel structure gives the attacker no advantage.

\paragraph{Effect of the Hankel structure.}
These estimates ignore the structure of $H$, and an attacker might exploit it. Its rows are shifts of one another, and the same sequence $a$ generates both $b$ and $u$, much as ML-KEM reuses its matrix $A$. We are not aware of an attack that improves on the unstructured estimates for this single-sample Hankel setting. However, the setting has received little study. Steinfeld et al.\ analyze lattice attacks on MP-LWE taking its Toeplitz structure into account~\cite{ssz2019practical}. Adapting that analysis to our parameters is the most important open question for this scheme.

\subsection{Decryption failure rate}\label{sec:dfr}
\begin{theorem}
The decryption failure rate satisfies $\mathrm{DFR} \le 2^{-134.1}$.
\end{theorem}
\begin{proof}
For each window coordinate $w$, the error $e_{\mathrm{dec}}[w]$ is a sum of 1{,}281 independent terms. Of these, 1{,}280 are products $X \cdot Y$ of independent $\CBD_2$ variables, and one is a $\CBD_2$ variable ($e_{2,w}$). Each term has mean $0$ and variance $1$. Hence $\sigma = \sqrt{1281} \approx 35.79$, and the threshold $512$ lies $14.31\sigma$ from the mean.

Evaluating the exact moment-generating functions of these distributions, the Chernoff bound gives $\Pr[e_{\mathrm{dec}}[w] \ge 512] \le 2^{-142.1}$. By symmetry, $\Pr[|e_{\mathrm{dec}}[w]| \ge 512] \le 2^{-141.1}$. A union bound over the $W = 2^7$ coordinates gives $\mathrm{DFR} \le 2^{-134.1}$.
\end{proof}

In 500 end-to-end IND-CCA exchanges (64{,}000 decoded coordinates), the largest observed $|e_{\mathrm{dec}}|$ was 131. The mean of $|e_{\mathrm{dec}}|$ was 28.4 and its standard deviation 21.4, close to the values of 28.5 and 21.6 predicted by a normal approximation. These observations are consistent with the bound; an experiment of this size cannot itself confirm a failure rate of $2^{-134}$.

%=====================================================================
\section{Implementation}\label{sec:impl}
%=====================================================================

We provide a Python reference implementation and an ARMv8 NEON implementation of the Hankel products and of decoding.

\paragraph{Accumulation without intermediate reductions.}
The forward product $\sum_{j<n} a_{i+j} s_j$ multiplies 16-bit sequence elements by 16-bit signed secret coefficients and accumulates the products in 32-bit lanes using \texttt{vmlal\_s16}. The largest possible magnitude is $640 \cdot 4095 \cdot 2 = 5{,}241{,}600$, far below $2^{31}$. Each output therefore needs only a single reduction (a mask with \texttt{0x0FFF}).

\paragraph{Branch-free decoding and selection.}
The 4-PAM decoder is applied to eight coordinates at a time:
\begin{lstlisting}[language=C, basicstyle=\small\ttfamily, frame=single]
uint16x8_t vw = vld1q_u16(in_w + i);
vw = vaddq_u16(vw, vdupq_n_u16(512));
vw = vandq_u16(vw, vdupq_n_u16(0x0FFF));
vw = vshrq_n_u16(vw, 10);
\end{lstlisting}
The ciphertext comparison and the implicit-rejection selection use bitwise select operations.

\paragraph{Performance and test vectors.}
The KEM is orchestrated in Python, with hashing, sampling and serialization in Python and NumPy; the NEON code computes the Hankel products and decoding. On an Apple M3 Max (Python 3.13, NumPy 2.3), the median latencies over 2{,}000 fresh key pairs and ciphertexts are 0.07~ms for key generation, 0.11~ms for encapsulation and 0.15~ms for decapsulation, and 0.34~ms for a full handshake. These figures include Python overhead and are not representative of a fully native implementation. We do not report timings for other schemes, because they were not measured on the same platform under the same conditions. Fifteen deterministic known-answer test vectors, generated with a custom SHAKE-based generator rather than the NIST DRBG, are reproduced exactly by the implementation.

%=====================================================================
\section{Limitations}
%=====================================================================
\begin{itemize}
  \item \emph{No reduction.} The instance lies outside the hypotheses of the MP-LWE reductions (Section~\ref{sec:mplwe}). Its security rests on concrete estimates for unstructured LWE, together with the assumption that the Hankel structure does not help an attacker.
  \item \emph{Unstructured estimates.} All attack costs in Section~\ref{sec:security} treat the instance as unstructured LWE. Attacks that exploit the Hankel structure have not been analyzed.
  \item \emph{No native implementation.} Only the arithmetic kernels are native; the reported timings include Python overhead.
  \item \emph{Quadratic arithmetic.} Without an NTT, multiplication costs $O(mn)$ operations, about $4.9 \times 10^5$ multiply-accumulates at $n = 640$. This is acceptable here, but it scales less favorably than ring-based schemes at higher security levels.
  \item \emph{Cryptanalytic maturity.} The single-sample Hankel setting has received little study. Independent cryptanalysis is needed before any use beyond research.
\end{itemize}

%=====================================================================
\section{Conclusion}
%=====================================================================
Windowed Hankel KEM publishes noisy middle products of a single pseudorandom sequence with a short secret. The result is a 2.5~KB exchange with a decryption failure rate below $2^{-134}$, and estimated lattice-attack costs at least 20 bits above those of Kyber-512 under the same cost models. The design gives up the MP-LWE reductions in exchange for compactness. Whether the Hankel structure can be exploited in this regime is the central open question.

\section*{Acknowledgments}
Portions of the analysis and the manuscript were prepared with the assistance of Claude (Anthropic), an AI model. The author takes full responsibility for the content.

\bibliographystyle{plain}
\bibliography{hankel_references}

\end{document}
